% !TeX program = xelatex
% analysis-of-alternatives.tex
% Document ENG-RPT-01 - Analysis of Alternatives: On-Premises ML Platform
%
% Build (from this folder, requires XeLaTeX; shared/intake-style.sty is a copy of the
% engineering_doc_templates version, kept in this repo so it builds on its own):
%   xelatex -interaction=nonstopmode analysis-of-alternatives.tex
%   xelatex -interaction=nonstopmode analysis-of-alternatives.tex   (rerun for TOC/page refs)
%
% Static report (no fillable fields). Shares branding, page layout, cover page
% and header/footer markings with the engineering intake template family via
% intake-style.sty. Companion document: ENG-RPT-02 Business Justification.
%
% Main file: preamble, cover page, document control and the section list. Each
% section is its own file in sections/ (edit those). Table column weights (Y{...})
% in each table must sum to the number of Y columns so the table fills \textwidth.

\documentclass[11pt]{article}

\usepackage{../../shared/intake-style}
\usepackage{../../shared/report-components}

\SetDocIdentity{Analysis of Alternatives: On-Premises ML Platform}{ENG-RPT-01}{v0.4-draft}

% Replace these placeholders once branding and markings are confirmed.
\SetCompanyName{[Engineering Center Name]}
\SetMarkingText{FOR OFFICIAL USE ONLY}

\SetDocPurpose{%
This report compares ways to provide a large-context, multi-user large
language model (LLM) capability that can be used with Controlled Unclassified
Information (CUI). It evaluates a proposed two-computer, on-premises design:
Computer~1, a network-attached storage (NAS) and private-cloud server that
hosts file shares, storage, software licensing and the LLM application tier;
and Computer~2, one or two machine-learning (ML) nodes that serve the model
and run fine-tuning, with the second node providing training capacity and
backup. The report gives price points, hardware, sizing logic, resilience and
growth options, risks and open items, so that a decision authority can choose
a configuration. It is a draft for later formal publication; all prices and
performance figures are planning estimates that must be replaced with vendor
quotes and an on-site benchmark before any purchase.%
}

\begin{document}

\CoverPage

\TodoAnchor{aoa-ctrl}
\ReportControlBlock{Draft}{v0.4}{2026-10-22}{\PH{Author}}{\PH{Intake Ref. No.}}{\PH{Recipients}}

\RevisionHistory{%
  v0.1 & \PH{Author} & 2026-10-14 & Initial AoA (large-budget tiers, single ML platform). \\
  v0.2 & \PH{Author} & 2026-10-16 & Re-based to a \$15K--\$55K ML-server budget. \\
  v0.3 & \PH{Author} & 2026-10-22 & Two-computer design (NAS + ML node[s]); dual-node training/backup option; multi-node growth path; LaTeX format. \\
  v0.4 & \PH{Author} & 2026-10-22 & Sponsor answers: $\geq$20 named / 12 simultaneous users; archive retrieval and training; Node~B as optional second serving node; Autodesk/Bluebeam hosting; Dell sourcing; cloud fallback options. \\
}

\begin{ReportNote}{Estimates, not quotes}
All prices, power, throughput and concurrency figures are planning-grade
estimates (typically $\pm$25--40\%) based on public street pricing and published
model characteristics. GPU, memory and disk prices are volatile. Items in
\PH{brackets} are placeholders that still need a value. This document contains
no CUI; verify that before adding project-specific detail.
\end{ReportNote}

% ===========================================================================
\IntakeSection{Items To Complete Before Release}

\Guidance{Every placeholder and open decision in this report is listed here and links to the
place it appears. Resolve each item, tick the Status box, and update the revision history
before the report is sent outside the team. ``Required'' items are missing information;
``Confirm'' items are assumptions that need a named person to agree to them.}\\[4pt]

{\footnotesize
\begin{xltabular}{\textwidth}{@{}c Y{1.5} Y{0.7} c c Y{0.8}@{}}
  \toprule
  \# & Item & Needed from & Priority & Status & Go to \\
  \midrule
  \endhead
  1 & Document identity and handling: author, intake reference number, recipients, center name and logo, marking text, version and date. & Document owner, security office & Required & $\square$ & \TodoLink{aoa-ctrl}{Control block} \\
  2 & Named-user growth beyond 20, workflow inventory, and the other planning assumptions (prompt sizes, speed targets). & Sponsor & Confirm & $\square$ & \TodoLink{aoa-assump}{Key assumptions} \\
  3 & Recovery targets (RTO/RPO) for LLM serving, the file share and the license services. & Sponsor, IT & Required & $\square$ & \TodoLink{aoa-assump}{Key assumptions} \\
  4 & Archive: size now, growth rate, file formats, folder and permission structure, and which documents to ingest first. & Sponsor, IT & Required & $\square$ & \TodoLink{aoa-storage}{Computer 1 needs} \\
  5 & Autodesk and Bluebeam: products, named-user or network licensing, internet sign-in, server components; other software to host. & IT, reseller, vendors & Required & $\square$ & \TodoLink{aoa-licenses}{Hosting Autodesk and Bluebeam} \\
  6 & Confirm that the license for each standard (for example NFPA~780) allows internal indexing and AI training. & Legal, publisher & Required & $\square$ & \TodoLink{aoa-nfpa}{Archive knowledge} \\
  7 & CUI categories in scope, export-controlled data, contract clauses, model-origin policy, SIEM log retention, and whether any cloud processing of CUI is allowed. & Security office, contracts & Required & $\square$ & \TodoLink{aoa-cui}{Compliance requirements} \\
  8 & Failover decision: is 3--8 users acceptable while Node~A is down, or must all 12 be served (Node~B with 2 GPUs, D4)? Will Node~B also serve routinely? & Sponsor & Required & $\square$ & \TodoLink{aoa-nodeb}{Node B modes} \\
  9 & Choose the cloud fallback option(s) to carry forward: none, non-CUI only, CUI-authorized government cloud, or warm standby. & Sponsor, security office & Required & $\square$ & \TodoLink{aoa-fallback}{Cloud fallback options} \\
  10 & Whether an offline backup copy is wanted in addition to disk redundancy; sponsor accepts the residual risk if not. & Sponsor & Confirm & $\square$ & \TodoLink{aoa-backup}{Redundancy} \\
  11 & Dell quotes for Computer~1 and each node configuration; which Dell models support the chosen GPUs; replace the 1.3 multiplier and the Dell-configured figures. & Dell account team & Required & $\square$ & \TodoLink{aoa-dell}{Dell-configured pricing} \\
  12 & Pilot benchmark: throughput and latency at 12 simultaneous users, long-context behavior, retrieval and fine-tuning quality and run times, ingestion rate; replace the estimates. & Project team & Required & $\square$ & \TodoLink{aoa-bench}{Sizing logic} \\
  13 & Facility readiness: circuits, cooling, noise, floor space, physical access; add 0.3~kW per GPU if Dell supplies 600~W GPUs. & Facilities & Required & $\square$ & \TodoLink{aoa-facility}{Facility} \\
  14 & Labor assumptions (hourly rate, upkeep hours, setup hours) and operating-cost estimates still accepted. & Sponsor, finance & Confirm & $\square$ & \TodoLink{aoa-tco}{Three-year cost} \\
  15 & Remaining open items: who operates both computers and how much time they have, server operating-system standard (Windows with Hyper-V or Linux with KVM), funding and procurement vehicle, approved marking text. & Sponsor & Required & $\square$ & \TodoLink{aoa-open}{Open items list} \\
  16 & Cross-check all figures against ENG-RPT-02 and update both documents together. & Document owner & Required & $\square$ & \TodoLink{aoa-tco}{Three-year cost} \\
  17 & No bracketed placeholders remain in the finished PDF; compiled twice so the contents and cross-references are current; revision history updated. & Document owner & Required & $\square$ & --- \\
  \bottomrule
\end{xltabular}}

% ===========================================================================
\ReportSection[sec:exec]{Executive Summary}

The organization needs an LLM capability that can process CUI with a large
context window, serve several people at once, and fit engineering workflows
(document analysis, retrieval over internal repositories, code and script
assistance, report drafting). The two tools currently provided lack the
functionality and context size for these workflows (see the companion Business
Justification).

This revision plans \textbf{two computers} rather than one:

\begin{itemize}
  \item \textbf{Computer~1 -- NAS and private-cloud server.} File share,
    bulk storage and backups, software license servers, a private cloud for
    specific software, and the LLM application tier (gateway, chat interface,
    retrieval, vector database, model and artifact registry, logging).
  \item \textbf{Computer~2 -- ML node(s).} Either a \textbf{single node} that
    serves the model, or \textbf{two nodes}: Node~A serves production use
    while Node~B runs fine-tuning and experiments at the same time, and can
    take over serving if Node~A fails.
\end{itemize}

The design is built so more nodes can be added later: every ML node is the same
class of machine, sits on the same storage network, and is placed behind the
same gateway.

\TableTitle[tbl:configs]{System configurations and planning price points (hardware only)}
{\footnotesize
\begin{xltabular}{\textwidth}{@{}l Y{1.062} l Y{0.692} Y{1.246}@{}}
  \toprule
  Config. & Composition & Hardware & Active users* & Training and resilience \\
  \midrule
  \endhead
  \textbf{S1} & Computer~1 + one 1-GPU node & \$35K (\$30--40K) & 3--8 & Training shares the GPU with serving (run off-hours). No failover. \\
  \textbf{S2} & Computer~1 + one 2-GPU node & \$44K (\$39--50K) & 8--20, or 3--8 with one GPU training & Split GPUs between serving and training. No failover. \\
  \textbf{D1} & Computer~1 + two 1-GPU nodes & \$55K (\$48--62K) & 3--8 & Dedicated training node; Node~B can serve at the same capacity if Node~A fails. \\
  \textbf{D3} & Computer~1 + one 2-GPU node (A) + one 1-GPU node (B) & \$64K (\$57--72K) & 8--20 & Dedicated training node; Node~B serves at reduced capacity (3--8) if Node~A fails. \\
  \textbf{S4} & Computer~1 + one 4-GPU node & \$65K (\$58--73K) & 25--50, or 3+1 GPU split & Single node; no failover. \\
  \textbf{D4} & Computer~1 + two 2-GPU nodes & \$73K (\$66--82K) & 8--20 (A); up to 16--40 with both serving & Dedicated training node; Node~B fails over at full capacity. \\
  \bottomrule
\end{xltabular}}
{\footnotesize * Simultaneous active users on the serving node(s). The stated need is at least 20 named users with about 12 at once, so S1 and D1 (1-GPU serving node) fall short. Hardware is at street prices; Dell-configured prices are about 1.15--1.4 times higher (Section~\ref{sec:dell}). Estimates to be confirmed by benchmark.}

\textbf{Confirmed requirements (v0.4).} At least 20 named users, with about 12 using
the model at the same time. The model must use the archived data on the NAS, both by
retrieval and by training on selected documents (for example NFPA~780). Training
alongside serving is the main use of the second node, which may also run a copy of
the serving model as a second node. The NAS will host Autodesk products and Bluebeam
\PH{and other software}. Hardware would be ordered from Dell, configured to
specification. Computer~1 gets its redundancy from disks inside the server rather
than from a second server.

\textbf{Preliminary recommendation.} Buy \textbf{D3}: one NAS and private-cloud server, a
2-GPU serving node (Node~A) and a 1-GPU training node (Node~B) that can also serve as a
second node, using 4-GPU-capable chassis so each node can grow by about \$9K per GPU at
street prices. Hardware is about \$64K at street prices and about \$83K configured by
Dell (planning multiplier 1.3, range 1.15--1.4; Section~\ref{sec:dell}). A 1-GPU serving
node is rated for 3--8 simultaneous users, below the stated 12, so S1 and D1 are not
recommended. If failover at full capacity is required, add a second GPU to Node~B (D4).
Alternative~2 (a compliant cloud LLM) remains a possible fallback; options are in
Section~\ref{sec:fallback}. See Section~\ref{sec:comparison} for the 3-year cost and
scoring, and Section~\ref{sec:open} for the decisions and information still needed.

\textbf{Cost note.} Computer~1 (about \$15K at street prices, plus its setup and upkeep)
is needed for the storage, licensing and private-cloud purposes whichever LLM option is
chosen, so the cost that differs between alternatives is the ML node(s).
Section~\ref{sec:tco} reports both the full-system and ML-only costs.

% ===========================================================================
\ReportSection[sec:req]{Problem Statement and Requirements}

\ReportSubsection{Operational needs}
\TodoAnchor{aoa-storage}

\textbf{Computer~1 (NAS and private cloud)}
\begin{itemize}
  \item Departmental file share (SMB, with NFS for the ML nodes), with access
    control lists, snapshots and quotas. Capacity target: \PH{N~TB now, growth
    \%/yr}.
  \item Bulk project and data storage; backup target for the ML nodes.
  \item License server hosting for \PH{named engineering software}.
  \item A private cloud for specific software: virtual machines and containers
    for \PH{named applications, e.g., version control, container registry,
    notebooks, wiki}.
  \item The LLM application tier: gateway and API, chat interface, retrieval
    service, vector database, model and artifact registry, centralized logging.
\end{itemize}

\textbf{Computer~2 (ML node[s])}
\begin{itemize}
  \item Serve an LLM to at least 20 named users, about 12 of them at the same
    time, with at least 128K-token context per request.
  \item Run fine-tuning and evaluation jobs (adapter-style fine-tuning of
    open-weight models, embedding and reranker adaptation, small custom ML
    models) \emph{at the same time} as serving, when two nodes are provisioned.
  \item Provide a backup path for serving if the primary node is unavailable.
  \item Provide an OpenAI-compatible API for scripts and pipelines, and
    retrieval over controlled repositories with per-group access control.
  \item Use the archived data on the NAS: answer questions from it with citations
    (retrieval), and train on selected documents, such as a licensed copy of NFPA~780,
    so the model handles that material well (Section~\ref{sec:archive}).
  \item Optionally load the same serving model on Node~B so it can act as a second
    serving node when it is not training (Section~\ref{sec:nodeb}).
\end{itemize}

\ReportSubsection{Compliance requirements (CUI)}
\TodoAnchor{aoa-cui}

\TableTitle{CUI-related requirements and their implications}
{\footnotesize
\begin{xltabular}{\textwidth}{@{}l Y{0.88} Y{1.12}@{}}
  \toprule
  Area & Driver & Implication \\
  \midrule
  \endhead
  Safeguarding & NIST SP 800-171 (Rev~2 or 3 per contract); 32~CFR Part~2002 & Access control, audit, encryption and boundary protection for every component, including the NAS. \\
  Contractual & DFARS 252.204-7012; CMMC Level~2 (if applicable) & Any cloud option must meet FedRAMP Moderate-equivalent and incident-reporting terms. \\
  Cryptography & FIPS 140-2/140-3 validated modules & FIPS-mode OS and encrypted disks. Confirm that the chosen hypervisor/NAS operating system supports this; an appliance NAS OS may not. \\
  Identity & CAC/PIV or enterprise MFA; role-based access & OIDC/SAML or Kerberos integration for file share, gateway and consoles. \\
  Audit & Who asked what; which documents were retrieved; who accessed which files & Logs forwarded to the SIEM with retention \PH{period}. \\
  Backup and retention & Contract and records rules & Backups stay inside the authorization boundary; encrypted; restore tested. \\
  Derived artifacts & Fine-tuned weights, embeddings and indexes built from CUI & Treated as CUI: stored in the access-controlled registry and covered by backup and sanitization rules. \\
  Supply chain & Model provenance; hardware sourcing & Approved model list; hash-verified offline import; TAA/Section~889-compliant hardware (see Section~\ref{sec:procure}). \\
  Export & ITAR/EAR data may be present & On-premises avoids foreign-hosted data; U.S.-person access control. \\
  Network & Air-gapped or enclave-only & Offline import of models and packages; internal mirrors. \\
  Authority & ATO / System Security Plan & One boundary should cover both computers; plan the authorization effort. \\
  \bottomrule
\end{xltabular}}

\ReportSubsection{Key assumptions (to validate)}
\TodoAnchor{aoa-assump}
\begin{enumerate}
  \item At least 20 named LLM users, with about 12 using the model at the same time
    (both stated by the sponsor). That is about 60\% simultaneous use, far above the
    10--20\% typical of office use, so the design is sized to 12 simultaneous users
    rather than to a ratio. Expected growth in named users: \PH{N}.
  \item Typical prompt 4--16K tokens (with retrieval); long jobs up to 128K.
    Typical response 300--1,500 tokens.
  \item Targets: time to first token under 3~s (short prompt) and under 20~s
    (128K prompt); at least 15--20 tokens/s per user while streaming.
  \item ``Training'' means fine-tuning and adaptation of existing open-weight
    models and small custom models. Pre-training a new LLM is out of scope at
    this budget.
  \item Failover target: \PH{RTO, e.g., manual switchover within 1~hour} for
    LLM serving. The file share and license services have their own recovery
    targets: \PH{RTO/RPO}.
  \item A controlled machine room or rack with 208/240~V circuits is available
    (GPU nodes are loud and hot; an open office is not suitable).
  \item Open-weight models are permitted after security review.
  \item Hardware is ordered from Dell, configured to specification. Prices here are street
    prices with a planning multiplier of 1.3 (range 1.15--1.4) until a quote is received.
  \item Storage need (best estimate, to confirm): the archive is about 20--40~TB now and
    grows 20--25\% per year, so Computer~1 needs about 40--70~TB usable within 3 years,
    including snapshots, vector indexes and model artifacts.
  \item The labor assumptions in Section~\ref{sec:tco} were accepted by the sponsor for now.
\end{enumerate}

% ===========================================================================
\ReportSection{Evaluation Criteria}

\TableTitle{Criteria and weights (weights revised in v0.3 to reflect the backup requirement)}
{\small
\begin{tabularx}{\textwidth}{@{}Y{0.667}Y{1.333}@{}}
  \toprule
  Criterion (weight) & Notes \\
  \midrule
  CUI compliance and data control (25\%) & Pass/fail gate plus degree of control. \\
  Capability: model quality and context (15\%) & Reasoning, coding, long-context accuracy. \\
  Concurrency and scalability (15\%) & Simultaneous users and growth path. \\
  Resilience and continuity (15\%) & Failover, backup, and impact of a single failure. \\
  Total cost of ownership, 3 years (15\%) & Hardware, setup and upkeep labor, operating cost. \\
  Time to capability (5\%) & Procurement lead time and authorization. \\
  Integration and customization (10\%) & API, retrieval, fine-tuning, workflow tooling. \\
  \bottomrule
\end{tabularx}}

% ===========================================================================
\ReportSection[sec:alts]{Alternatives Considered}

\TableTitle{Alternatives and disposition}
{\footnotesize
\begin{xltabular}{\textwidth}{@{}l Y{0.667} Y{1.333}@{}}
  \toprule
  \# & Alternative & Disposition \\
  \midrule
  \endhead
  0 & \textbf{Status quo}: two existing in-house tools & Baseline. No new cost; capability and context gaps persist; limited API and retrieval. Does not provide the storage, licensing or private-cloud capability. \\
  1 & \textbf{Commercial public SaaS} & \textbf{Eliminated.} Not suitable for CUI. \\
  2 & \textbf{Compliant cloud LLM} (e.g., Azure OpenAI in Azure Government, AWS Bedrock in GovCloud) & Viable for the LLM only. Fast, strong models, no GPU hardware. Cost scales with long-context use; CUI leaves the site; contract and boundary review needed. Estimated \$0.03M--\$0.4M over 3 years for 20--30 users (usage-driven; verify). \textbf{Computer~1 is still required} for storage and licensing. \\
  3 & \textbf{On-premises design (Computer~1 + ML node[s])} & Detailed in Section~\ref{sec:hw}. Two families: \textbf{single-node} (S1, S2, S4) and \textbf{dual-node} (D1, D3, D4). \\
  4 & \textbf{Hybrid}: on-premises for CUI, cloud for non-CUI & Possible later; needs a routing policy and adds mis-routing risk. \\
  \bottomrule
\end{xltabular}}

\textbf{Below the smallest ML node (not recommended for production).} Small
unified-memory systems (for example NVIDIA DGX Spark-class or AMD Strix Halo
mini PCs, about \$2--5K) can run large models for one developer, but their
memory bandwidth is roughly 6--7 times lower than a 96~GB workstation GPU, so
they do not serve several simultaneous users well.

\ReportSubsection[sec:fallback]{Cloud fallback options}
\TodoAnchor{aoa-fallback}

Whether a cloud service should back up the on-premises LLM is undecided; the options
below are kept open for refinement.

\TableTitle{Fallback options if the on-premises LLM is unavailable}
{\footnotesize
\begin{xltabular}{\textwidth}{@{}l Y{1.0} Y{1.1} Y{0.9}@{}}
  \toprule
  Option & Description & CUI and network considerations & Cost \\
  \midrule
  \endhead
  F0 & No cloud fallback. Users use Tool A or Tool B, or wait. & Nothing leaves the boundary. & None. \\
  F1 & Commercial cloud LLM for non-CUI work only. & No CUI may be entered; needs a user policy and an approved internet path outside the enclave. & Low, pay per use. \\
  F2 & Government-cloud LLM authorized for CUI (for example Azure OpenAI in Azure Government or Bedrock in GovCloud), reached over an approved connection. & Needs contract flow-down, an authorization decision, and a network path that an air-gapped design would not have. Data leaves the site. & Pay per use; setup effort for the account, connection and approval; usually no cost while idle. \\
  F3 & Warm standby: the same open-weight stack, deployed from configuration code onto rented compliant GPU instances only when needed. & Same boundary questions as F2; CUI copies of indexes would be needed there. & Hourly instance cost while running; setup effort. \\
  \bottomrule
\end{xltabular}}
F0 is the baseline in this report. F2 is the most capable and the most demanding to
authorize. \PH{Select the option(s) to carry forward.}

% ===========================================================================
\ReportSection[sec:concept]{Proposed System Concept}

\ReportSubsection{Two-computer architecture}

\FigureTitle[fig:arch]{Logical architecture: Computer~1 (NAS and private cloud) and Computer~2 (ML nodes)}
\begin{center}
\begin{tikzpicture}[
  svc/.style={draw, rounded corners=1.5pt, font=\scriptsize, align=center,
              minimum height=0.8cm, minimum width=3.4cm, inner sep=2pt},
  host/.style={draw, rounded corners=3pt, inner sep=7pt},
  lbl/.style={font=\small\itshape},
  arr/.style={-{Stealth[length=2mm]}, thick},
  darr/.style={-{Stealth[length=2mm]}, thick, dashed},
  bi/.style={{Stealth[length=2mm]}-{Stealth[length=2mm]}, thick}
]
\node[svc, minimum width=2.4cm] (users) at (0,-1) {Users:\\browser, IDE,\\scripts};
\node[svc] (gw)  at (5.2,-1) {LLM gateway, chat UI,\\retrieval service};
\node[svc] (vdb) at (5.2,-2.2) {Vector database,\\model and artifact registry};
\node[svc] (fs)  at (5.2,-3.4) {File share (SMB/NFS),\\snapshots, backups};
\node[svc] (lic) at (5.2,-4.6) {License servers,\\private-cloud VMs};
\begin{scope}[on background layer]
  \node[host, fit=(gw)(vdb)(fs)(lic), label={[lbl]above:Computer 1: NAS / private cloud}] (c1) {};
\end{scope}
\node[svc, minimum width=4.2cm, minimum height=1.3cm] (na) at (12.7,-1) {ML Node A (primary)\\serving: inference,\\embeddings, reranker};
\node[svc, minimum width=4.2cm, minimum height=1.3cm] (nb) at (12.7,-3.4) {ML Node B (secondary)\\fine-tuning, experiments;\\standby for Node A};
\node[lbl, above=0.15cm of na] {Computer 2: ML nodes};
\draw[arr] (users) -- (gw);
\draw[arr] (gw.east) -- node[above, font=\scriptsize]{inference API} (na.west);
\draw[darr] (gw.east) ++(0.5,0) |- node[pos=0.25, right, font=\scriptsize]{failover} ([yshift=0.35cm]nb.west);
\draw[bi] (fs.east) -- node[below, font=\scriptsize]{datasets, checkpoints, backups} (nb.west |- fs.east);
\end{tikzpicture}
\end{center}
{\footnotesize Every node mounts the NAS over the storage network and keeps its own
local copy of the approved models, so serving continues if Computer~1 is
briefly unavailable. A second gateway instance can run on Node~B (Section~\ref{sec:resil}).}

\ReportSubsection{Computer 1: roles}
\begin{itemize}
  \item \textbf{File and data services:} SMB share with Active Directory or
    Kerberos authentication, NFS exports to the ML nodes, snapshot schedule,
    quotas and an encrypted backup target.
  \item \textbf{License servers:} hosted in virtual machines or on the host.
    \emph{Check each vendor's terms}: some license managers are bound to a
    hardware or MAC address and may restrict virtualization.
  \item \textbf{Private cloud:} a hypervisor (for example Proxmox or KVM on a
    FIPS-capable Linux) running the specific software \PH{list}, each in its
    own virtual machine or container, with snapshots and backup.
  \item \textbf{LLM application tier:} gateway, chat interface, retrieval
    service, vector database, model and artifact registry, observability.
\end{itemize}

\ReportSubsection[sec:licenses]{Hosting Autodesk and Bluebeam}
\TodoAnchor{aoa-licenses}

The stated software to host is Autodesk products and Bluebeam. These need to be checked
with the vendors or reseller before the Computer~1 design is frozen; the points below are
the questions to ask, not confirmed facts.

\TableTitle{Autodesk and Bluebeam: what to verify}
{\footnotesize
\begin{xltabular}{\textwidth}{@{}Y{0.636} Y{1.364} Y{1.0}@{}}
  \toprule
  Topic & Question & Why it matters \\
  \midrule
  \endhead
  Licensing model & Are the licenses named-user subscriptions that sign in through the vendor's cloud, or network (floating) licenses served from a local license manager? (Autodesk and Bluebeam have largely moved to named-user subscriptions; legacy agreements may still use a local license manager.) & If named-user, a local license server may not be needed; if network, the license manager runs on Computer~1. \\
  Internet sign-in & Does the product need periodic internet check-in or sign-in? & An enclave or air-gapped network may not allow it. \\
  Data and collaboration services & Which services are needed (for example Autodesk Vault for data management, Bluebeam collaboration)? Is a self-hosted option available and authorized for CUI? & Cloud versions sit outside the CUI boundary unless authorized. \\
  Server platform & Do the server components require Windows Server and SQL Server? & Adds Windows Server, database and client-access licensing, and may favor a Hyper-V host. \\
  Virtualization and hardware binding & Does the vendor support running the license manager in a virtual machine, and is it bound to a host ID or MAC address? & May require a fixed virtual MAC or a dedicated host. \\
  \bottomrule
\end{xltabular}}

\ReportSubsection{Computer 2: ML nodes and workloads}

\TableTitle{How each node is used}
{\small
\begin{tabularx}{\textwidth}{@{}Y{0.778}Y{1.111}Y{1.111}@{}}
  \toprule
  & Single-node configs & Dual-node configs \\
  \midrule
  Serving & Node serves all users. & Node~A serves; Node~B is warm standby. \\
  Fine-tuning & Time-shared with serving (off-hours) or on a separate GPU if the node has two or more. & Runs on Node~B concurrently with serving on Node~A. \\
  Experiments and evaluation & Same as fine-tuning. & Node~B; does not disturb production latency. \\
  Failure of the node & Service down until repaired or the cloud fallback is used. & Gateway routes to Node~B; training pauses until Node~A returns. \\
  Maintenance (patching, driver or model upgrades) & Planned downtime. & Rolling: patch Node~B first, then switch and patch Node~A. \\
  \bottomrule
\end{tabularx}}

A dual-node design is worth its extra cost mainly for three reasons: (1)
training cannot degrade user-facing latency, (2) there is a failover path, and
(3) upgrades can be tested on Node~B before production sees them. A single
node with two or more GPUs can also train and serve at the same time by
dedicating GPUs to each, but it shares one failure domain and one power supply.

\ReportSubsection[sec:archive]{Archive knowledge: retrieval and fine-tuning}

Archived project data on the NAS should be usable by the model. Two techniques do
different jobs, and the design uses both.

\TableTitle{Retrieval versus fine-tuning for archive knowledge}
{\footnotesize
\begin{xltabular}{\textwidth}{@{}Y{0.632} Y{1.158} Y{1.053} Y{1.157}@{}}
  \toprule
  Technique & What it does & Strengths & Limits \\
  \midrule
  \endhead
  Retrieval (index the documents) & Splits documents into passages, stores them in a searchable index, and gives the model the best-matching passages with each question. & Quotes exact wording and cites document, clause and page; new files are usable as soon as they are indexed; follows NAS permissions; no training time. & Answer quality depends on finding the right passages; limited by the context size per question. \\
  Fine-tuning (train an adapter) & Adjusts the model using example questions and answers or documents. & Learns vocabulary, formats and how to do tasks (for example draft an inspection checklist in the house format); shorter prompts. & Does not reliably memorize a document and can state wrong clause numbers with confidence; content is hard to update or remove; must be re-tested after each change; the resulting weights are CUI. \\
  \bottomrule
\end{xltabular}}

Giving the model a PDF and training on it will not by itself make it a dependable
expert. For a standard such as NFPA~780, where exact wording matters, \textbf{retrieval
with citations is the primary mechanism}; fine-tuning on reviewed examples teaches the
model how to use and reason about the material. A fine-tuned adapter is deployed only
if it measurably beats retrieval alone.

\textbf{Example workflow: NFPA~780.}
\begin{enumerate}
  \item Place the licensed PDF in a controlled library folder on the NAS with the intended
    access groups.
  \item An ingestion job (run on Node~B or on CPU) extracts text and tables, applies optical
    character recognition to scanned pages, splits by section and clause, creates
    embeddings, and stores them in the vector database with edition and access metadata.
  \item Users ask questions; the model answers with citations to clause and page.
  \item Optional training: example questions, answers and tasks are generated from the
    document and reviewed by an engineer; an adapter is trained on Node~B; it is evaluated
    against a held-out set of questions with known answers and compared with retrieval alone.
  \item When a new edition is issued, re-ingest it with the edition recorded so that old and
    new requirements are not mixed.
\end{enumerate}

\TodoAnchor{aoa-nfpa}
\textbf{Points to resolve.}
\begin{itemize}
  \item \textbf{Formats.} Text-based and scanned documents work. CAD and drawing files (for
    example DWG or Revit) are not text; in the base plan only titles, annotations and
    metadata can be extracted. Reading drawings visually needs a vision model that is not
    budgeted here.
  \item \textbf{Publisher licenses.} Standards such as NFPA~780 are copyrighted. Confirm that
    the license allows internal indexing and AI training.
  \item \textbf{Access control.} Retrieval must respect NAS permissions. A model trained on a
    restricted folder must not be available to people without access, so train separate
    adapters per access group, or train only on material open to all users of that model.
  \item \textbf{Derived data.} Indexes, embeddings and adapters built from CUI are treated as
    CUI and covered by backup and sanitization rules.
  \item \textbf{Effort.} Ingestion, OCR and training run on Node~B and store results on
    Computer~1; archive size and formats (open item) set the effort.
\end{itemize}

\ReportSubsection[sec:nodeb]{Node B operating modes}
\TodoAnchor{aoa-nodeb}

Training alongside serving is the main use of Node~B. Because it is the same class of
machine as Node~A, it can also run a copy of the serving model. This is an option within
the design, not a fixed decision.

\TableTitle{Node B modes}
{\footnotesize
\begin{xltabular}{\textwidth}{@{}Y{0.656} Y{1.031} Y{1.313}@{}}
  \toprule
  Mode & When & Effect on capacity \\
  \midrule
  \endhead
  1. Training (default) & Fine-tuning, document ingestion, OCR and evaluation jobs run on Node~B. & Node~A serves users alone, so capacity is Node~A's. \\
  2. Second serving node & No training job is running; Node~B loads the same model and the gateway shares the load. & Adds about 3--8 simultaneous users (1 GPU) or 8--20 (2 GPUs). Training jobs take priority: the gateway drains Node~B sessions before a job starts. \\
  3. Failover & Node~A is unavailable. & Node~B serves alone at its own capacity (3--8 users with 1 GPU); training is paused. \\
  \bottomrule
\end{xltabular}}

A 1-GPU Node~B cannot train and serve at the same time: one 96~GB GPU holds either the
serving model with its cache (about 63~GB plus cache) or a training job (about 20--90~GB).
A 2-GPU Node~B (D4) can do one of each. Training jobs should save checkpoints often so
they can be paused for failover or for serving. A job queue (for example Slurm or
Kubernetes with priority classes) enforces the policy; a typical policy is training at
night and on weekends, and serving on Node~B during working hours when it is idle.

\ReportSubsection{Software stack (self-hosted, no outbound connectivity)}

\TableTitle{Software layers}
{\footnotesize
\begin{xltabular}{\textwidth}{@{}l Y{1.231} Y{0.769}@{}}
  \toprule
  Layer & Candidates & Purpose \\
  \midrule
  \endhead
  OS and virtualization & RHEL, Rocky or Ubuntu Pro in FIPS mode with STIG hardening; KVM/Proxmox; Docker or k3s & Hardened base; private cloud. \\
  Storage & ZFS (RAID-Z2) or equivalent, SMB/NFS, snapshot and replication tooling & File share, backups, datasets. \\
  Inference engine & vLLM or SGLang (TensorRT-LLM optional) & Continuous batching, paged KV cache, prefix caching. \\
  Training tools & PEFT/LoRA frameworks, experiment tracking (e.g., MLflow) & Adapter fine-tuning, evaluation. \\
  Gateway & LiteLLM or NGINX with OIDC, health checks, optional keepalived virtual IP & OpenAI-compatible API, quotas, failover routing. \\
  User interface & Open WebUI or similar & Chat, uploads, shared prompts. \\
  Retrieval & Open embedding model and reranker; pgvector or Qdrant & Cited answers with access-control enforcement. \\
  Document ingestion & PDF and Office text extraction, OCR for scanned pages, embedding model & Turn NAS archive content into searchable passages. \\
  Observability & Prometheus, Grafana, DCGM exporter, log forwarding to SIEM & Utilization, latency, audit. \\
  Configuration management & Ansible or similar, stored in the NAS git server & Identical node builds; fast rebuild. \\
  \bottomrule
\end{xltabular}}

\ReportSubsection{Model selection approach}
Use a portfolio sized to the node:

\TableTitle{Model roles and memory}
{\small
\begin{xltabular}{\textwidth}{@{}Y{0.844}Y{1.5}Y{0.656}@{}}
  \toprule
  Role & Candidate classes (verify versions, licenses, provenance) & Weights memory \\
  \midrule
  \endhead
  General assistant & OpenAI gpt-oss-120b (mixture of experts, Apache 2.0, about 5B active parameters) & about 63~GB (4-bit) \\
  Mid-size general or coding & Llama 3.x 70B, Gemma 27B, Mistral Small/Medium class, NVIDIA Nemotron & 15--40~GB (4-bit) to 70~GB (FP8) \\
  Large mixture of experts (4-GPU node only) & Llama 4 Maverick-class (about 400B total) at 4-bit & about 200--220~GB \\
  Embeddings and reranker & Small open models & 1--4~GB \\
  \bottomrule
\end{xltabular}}

\textbf{Provenance policy.} This report assumes U.S. or allied-origin models
only; many organizations prohibit models from foreign-adversary developers on
government-related systems regardless of license. Download on an internet-connected
system, scan, hash-verify, store as \texttt{safetensors} (no pickle files), and
import through the approved transfer process. Legal review of each license is
required. Mixture-of-experts models are well suited to this budget: few active
parameters per token mean one 96~GB GPU can serve several users quickly, and
the cache memory per token is small.

\ReportSubsection{How multiple users share the system}
\begin{enumerate}
  \item \textbf{Continuous batching:} the inference engine interleaves many users'
    requests token by token, so one loaded model serves everyone at once.
  \item \textbf{Paged KV cache:} each active conversation uses GPU memory in
    proportion to its context length; memory is usually the limit for long
    context (Section~\ref{sec:sizing}).
  \item \textbf{Prefix caching:} shared system prompts and repeated documents are cached.
  \item \textbf{Quotas and priority:} the gateway enforces per-user and per-team token
    budgets and prioritizes interactive use over batch jobs.
  \item \textbf{Isolation:} conversations and retrieval indexes are partitioned by
    user and group; logging is per user.
  \item \textbf{Batch queue:} large offline jobs run off-peak.
\end{enumerate}

\ReportSubsection[sec:resil]{Resilience, backup and multi-node growth}

\TableTitle{Failure modes and responses}
{\footnotesize
\begin{xltabular}{\textwidth}{@{}Y{0.8}Y{1.0}Y{1.0}Y{1.2}@{}}
  \toprule
  Failure & Single-node configs & Dual-node configs & Mitigation \\
  \midrule
  \endhead
  ML node A down & LLM unavailable. & Gateway fails over to Node~B; training paused. & Health checks; models staged locally on both nodes; documented switchover. \\
  ML node B down & Not applicable. & No user impact; no failover or training until repaired. & Spare parts; warranty. \\
  One GPU fails & Node continues with fewer GPUs if the model fits; otherwise down. & Same, plus failover. & Keep one cold-spare GPU if the budget allows. \\
  Computer 1 (NAS) down & File share, license servers, private cloud and the gateway/retrieval tier are down. & Same for NAS services; LLM serving continues if a second gateway runs on Node~B (retrieval corpus unavailable). & Redundant PSUs, RAID-Z2, UPS; VM backups; tested restore; optional replica (Section~\ref{sec:hw}). \\
  Disk failure on NAS & Array degraded but online. & Same. & RAID-Z2 tolerates two disks; hot spare or fast replacement. \\
  Power loss & All down. & All down. & UPS with graceful shutdown; generator if mission-critical. \\
  Bad model or software update & Rollback only. & Test on Node~B first. & Versioned registry; configuration as code. \\
  Site loss (fire, flood) & Data loss unless backed up elsewhere. & Same. & Encrypted offline or off-site copy of backups within the CUI boundary. \\
  \bottomrule
\end{xltabular}}

\textbf{Growth path (multi-node).}
\begin{itemize}
  \item \textbf{Same build, more nodes.} Every ML node uses the same chassis class,
    operating-system image and configuration, so a third or fourth node is another
    copy. The storage switch has spare 25~GbE ports, and the gateway load-balances
    across any number of nodes.
  \item \textbf{Scale out by replicas.} More nodes serve more simultaneous users
    (each node holds a full copy of the model) or run more training jobs in
    parallel. Add GPUs inside a node (about \$9K each) before adding a node.
  \item \textbf{Limits.} These PCIe-based workstation GPUs do not support
    efficient multi-node tensor parallelism, so a single model replica must fit
    inside one node (up to 384~GB on a 4-GPU node). A single training job spanning
    several nodes would need an RDMA-capable 100--200~GbE fabric (additional
    network cards and switch, roughly \$7--15K for two nodes; estimate); this is
    not part of the base plan.
  \item \textbf{Computer 1 redundancy.} A second NAS/private-cloud server (about
    \$12--18K) with replicated storage and virtual machines is the next step if
    the license servers or file share become critical; see Section~\ref{sec:open}.
\end{itemize}

% ===========================================================================
\ReportSection[sec:components]{System Components and Service Expectations}

\ReportSubsection{Plain-language overview}

The system lets many people ask questions of a large language model (LLM) from their
own computers, with the model running entirely on equipment inside the organization.
A person types a question or uploads a document in a web page or a script. A gateway
checks who they are, a retrieval service finds relevant passages in the archived files
on the NAS, and the model, running on a GPU server, writes an answer that cites those
passages. Many people can do this at once because the model software handles their
requests together rather than one after another. A second GPU server is used mainly to
train the model on organization-specific material without slowing anyone down, and it
can also take over serving when needed. The first computer holds the files, backups,
license servers and the supporting software, and nothing leaves the organization's
boundary.

\ReportSubsection{What each piece is and does}

\TableTitle{System components}
{\footnotesize
\begin{xltabular}{\textwidth}{@{}Y{0.8} Y{1.3} Y{1.2} Y{0.7}@{}}
  \toprule
  Component & What it is & What it does for users & Runs on \\
  \midrule
  \endhead
  ML node (GPU server) & A server holding one or more 96~GB graphics processing units (GPUs). The GPU memory is where the model lives while it works. & Provides the computing power to answer questions quickly and to train the model. & Computer~2 \\
  Language model (LLM) & A large file of learned numbers (weights) that predicts text; ``open-weight'' means the organization may run it locally. & Reads questions and documents and writes answers, summaries, code and drafts. & ML node GPUs \\
  Inference server & Software (for example vLLM) that loads the model and serves requests. & Lets many users share one loaded model, streams answers as they are written, and keeps recent context in memory to save time. & ML node \\
  Embedding model and reranker & Small models that turn text into numbers that capture meaning, and that reorder search results by relevance. & Make search work by meaning rather than exact words. & ML node or CPU \\
  Vector database & A database of those numbers for every document passage. & Finds the passages most relevant to a question in milliseconds. & Computer~1 \\
  Ingestion pipeline & Programs that read files on the NAS, extract text (with OCR for scans), split it into passages, and index it. & Make archived documents and standards searchable; re-run when files change. & Node~B or CPU; results stored on Computer~1 \\
  Retrieval service & The logic that takes a question, finds passages, builds the prompt, calls the model, and returns an answer with citations. & Gives answers tied to documents, clauses and pages, and applies each user's access rights. & Computer~1 \\
  LLM gateway & A single front door with an OpenAI-compatible interface. & Signs users in, applies quotas and priorities, routes to Node~A or B, fails over, and logs each request. & Computer~1 \\
  Chat interface & A web page for chatting, uploading files and reusing prompts. & The everyday way most people use the system. & Computer~1 \\
  Training environment & Job queue, fine-tuning tools, experiment tracking and evaluation tests. & Adapts the model to organization vocabulary and formats; tests whether it helped. & Node~B \\
  Model and artifact registry & A versioned, access-controlled store of approved models, adapters and indexes with checksums. & Ensures only approved, verified models run and allows rollback. & Computer~1 \\
  Identity provider & The directory and sign-in service (for example Active Directory with CAC/PIV or multi-factor). & Controls who can use the system and which documents they can retrieve. & Existing or Computer~1 \\
  Logging and monitoring & Collection of audit logs and performance measurements. & Shows who asked what, supports compliance evidence, and warns of overload or failures. & Computer~1, forwarded to the SIEM \\
  File share (NAS) & Shared storage with permissions, snapshots and quotas, served over SMB and NFS. & Holds project data, the archive, datasets, checkpoints and backups. & Computer~1 \\
  License servers & Programs that hand out floating software licenses on the network. & Allow engineering software to run; needed only for network-type licenses. & Computer~1 (VM) \\
  Private cloud (hypervisor) & Virtualization software that runs several isolated virtual machines on one server. & Hosts Autodesk, Bluebeam and other specific software in separate, snapshot-protected machines. & Computer~1 \\
  Storage network and firewall & High-speed switch (25~GbE) and a firewall around the system. & Move data between computers quickly; keep the system inside the boundary. & Both \\
  \bottomrule
\end{xltabular}}

\ReportSubsection{Expected service levels (planning targets)}

These are targets used for sizing and for acceptance testing in the pilot. They are not
guarantees until measured on the real hardware with real prompts.

\TableTitle{Service expectations}
{\footnotesize
\begin{xltabular}{\textwidth}{@{}Y{1.0} Y{0.9} Y{1.1}@{}}
  \toprule
  Measure & Target & Notes \\
  \midrule
  \endhead
  Named users & At least 20 (stated) & Growth: \PH{expected named users in 3 years}. \\
  Simultaneous users generating answers & About 12 (stated) & D3 Node~A is rated 8--20; D4 or Node~B serving adds headroom. \\
  Context size per request & At least 128K tokens (about 300 pages of text) & Quality of recall falls with length; retrieval is preferred for very large archives. \\
  Time to first word & Under 3~s for prompts up to 16K tokens; under 20~s for 128K & Measured at the gateway under load. \\
  Streaming speed per user & At least 15--20 tokens/s (faster than reading speed) & Falls if many long prompts arrive together. \\
  Answers from the archive & Cite document, clause and page & An answer without a citation should be treated as unsupported. \\
  Availability in duty hours & Target \PH{98\%} & Single-node configurations have no ML failover; Computer~1 has no server redundancy. \\
  Recovery after failure & Target \PH{RTO} for LLM serving; \PH{RTO/RPO} for files and licenses & Dual-node: switch to Node~B; test before use and after major changes. \\
  Model refresh & Quarterly review of newer approved models & Tested on Node~B before production. \\
  Fine-tuning turnaround & Hours to a few days per adapter on one GPU & Measure in the pilot; depends on data size. \\
  Archive ingestion rate & \PH{documents per hour, measure in pilot} & OCR of scanned pages is the slowest step. \\
  Data handling & All CUI stays inside the authorization boundary; per-user audit logs & Derived data (indexes, adapters) is treated as CUI. \\
  \bottomrule
\end{xltabular}}

\ReportSubsection{What the system will and will not do}

\textbf{Expected to do well}
\begin{itemize}
  \item Answer questions from approved documents and cite the source passage.
  \item Summarize, compare and extract data from long documents, up to the context limit.
  \item Draft reports, procedures, emails, scripts and code for an engineer to review.
  \item Process large batches of documents overnight.
  \item Learn organization vocabulary, formats and task habits through fine-tuning.
  \item Serve about a dozen people at the same time with acceptable speed.
\end{itemize}

\textbf{Limits to set expectations}
\begin{itemize}
  \item It can be wrong. Every answer used in a deliverable needs engineer review; the
    system does not replace professional judgment or stamping.
  \item Training on a document does not reliably make the model remember it exactly;
    retrieval with citations is how exact wording is obtained.
  \item In the base plan it cannot read drawings or photographs as images; it reads
    text, tables and drawing metadata only.
  \item It does not browse the internet or call outside services.
  \item Pre-training a new model from scratch is out of scope at this budget.
  \item The best open-weight models can trail the best proprietary models on the hardest
    reasoning tasks.
  \item Capacity is finite: at peak, requests queue, and failure of the serving node
    reduces service until failover or repair.
\end{itemize}

% ===========================================================================
\ReportSection[sec:hw]{Hardware, Price Points and Configurations}

\ReportSubsection[sec:nas]{Computer 1: NAS and private-cloud server}

\TableTitle{Computer 1 reference bill of materials (planning estimates)}
{\footnotesize
\begin{xltabular}{\textwidth}{@{}Y{0.621}Y{1.379}r@{}}
  \toprule
  Component & Specification & Est. cost \\
  \midrule
  \endhead
  Chassis and platform & 2U/4U rack, 12 or more 3.5-inch bays, one 16--24-core server CPU, TPM~2.0, BMC, redundant power supplies & \$3--4K \\
  Memory & 256~GB ECC & \$1.5--2.5K \\
  System and VM storage & 2 $\times$ 1.9~TB NVMe (mirror, OS and hypervisor); 2 $\times$ 3.8~TB NVMe (mirror, VMs, vector database) & \$1--1.5K \\
  Bulk storage & 8 $\times$ 20~TB enterprise HDD in RAID-Z2 or RAID~6 (about 100~TiB usable; plan at most 80\% fill); 4 bays free for growth & \$3--3.6K \\
  Controller and network & SAS HBA; dual-port 25~GbE & \$0.5--0.8K \\
  UPS and rack & 3~kVA UPS shared with the ML nodes & \$1--2K \\
  OS and hypervisor subscription & FIPS-capable Linux and virtualization support & \$0.3--1K \\
  Switch, firewall, cabling & 25~GbE switch with spare ports, firewall, management network & \$1.5--2.5K \\
  \midrule
  \textbf{Computer 1 total} & & \textbf{\$12--18K (about \$15K)} \\
  \bottomrule
\end{xltabular}}

\textbf{Storage sizing (best estimate, to confirm).} The archive is assumed to be about
20--40~TB now, growing 20--25\% per year, so about 40--70~TB usable is needed within 3
years including snapshots, vector indexes and model artifacts. The 8 $\times$ 20~TB array
(about 100~TiB usable, filled to at most 80\%) covers this with room to spare, and 4
bays remain free.

\TodoAnchor{aoa-backup}
\textbf{Redundancy.} Following the sponsor's direction, Computer~1 is protected by its
disks rather than by a second server: RAID-Z2 (tolerates two disk failures), a hot spare
(add one 20~TB drive, about \$0.4K), redundant power supplies and a UPS, plus snapshots.
This covers disk failure and accidental deletion. It does not cover loss of the server,
ransomware or a site loss, because the data and the snapshots share one machine; an
encrypted offline backup set rotated off-site (about \$2--4K) closes that gap and is
recommended but optional.

Likely additions not included in the table: Windows Server, database and client-access
licensing for Autodesk and Bluebeam server components (about \$2--8K; estimate, depends on
the products chosen; Section~\ref{sec:licenses}).

\ReportSubsection{Computer 2: ML node building blocks}

\TableTitle{ML node options (each in a 4-GPU-capable chassis so GPUs can be added later)}
{\footnotesize
\begin{xltabular}{\textwidth}{@{}l Y{0.833} Y{1.167} r@{}}
  \toprule
  Node & GPUs & Host and notes & Est. cost \\
  \midrule
  \endhead
  \textbf{N1} & 1 $\times$ RTX PRO 6000 Blackwell 96~GB (Max-Q, 300~W), \$8--10K & 16--32-core CPU, 256~GB RAM, 2 $\times$ 3.8~TB NVMe, dual 1600--2000~W PSUs, 25~GbE: \$9--11K; misc.\ \$1K & \$18--22K \\
  \textbf{N2} & 2 GPUs (192~GB total), \$16--20K & Same class of host: \$10--11K; misc.\ \$1K & \$27--32K \\
  \textbf{N4} & 4 GPUs (384~GB total), \$32--40K & Host with 512~GB RAM, larger PSUs: \$12--13K; misc.\ \$2K & \$46--55K \\
  \bottomrule
\end{xltabular}}

Adding one GPU to an existing node costs about \$8--10K. Government-channel
(TAA-compliant OEM) pricing is often 15--40\% higher than the street prices used
here (Section~\ref{sec:procure}).

\ReportSubsection{System configurations and price ladder}

\TableTitle{Configurations (Computer~1 plus ML nodes)}
{\footnotesize
\begin{xltabular}{\textwidth}{@{}l Y{0.8} r r Y{1.2}@{}}
  \toprule
  Config. & ML node(s) & Hardware total & GPU memory & Named users \\
  \midrule
  \endhead
  \textbf{S1} & N1 & \$35K (\$30--40K) & 96~GB & 10--30 \\
  \textbf{S2} & N2 & \$44K (\$39--50K) & 192~GB & 25--75 \\
  \textbf{D1} & N1 + N1 (A + B) & \$55K (\$48--62K) & 2 $\times$ 96~GB & 10--30 \\
  \textbf{D3} & N2 (A) + N1 (B) & \$64K (\$57--72K) & 192 + 96~GB & 30--100 \\
  \textbf{S4} & N4 & \$65K (\$58--73K) & 384~GB & 75--250 \\
  \textbf{D4} & N2 + N2 (A + B) & \$73K (\$66--82K) & 2 $\times$ 192~GB & 30--100 (A); up to 60--200 with both serving \\
  \bottomrule
\end{xltabular}}

\TableTitle{Upgrade paths}
{\small
\begin{tabularx}{\textwidth}{@{}Y{1.0}Y{0.6}Y{1.4}@{}}
  \toprule
  Step & Added cost & Result \\
  \midrule
  S1 $\to$ S2 (add one GPU) & about \$9K & Doubles memory; train and serve at once on separate GPUs. \\
  S1 $\to$ D1 (add node B) & about \$20K & Dedicated training node and failover path. \\
  D1 $\to$ D3 (add GPU to node A) & about \$9K & Roughly doubles serving capacity; failover is at reduced capacity. \\
  D3 $\to$ D4 (add GPU to node B) & about \$9K & Failover at full capacity; both nodes can serve when not training. \\
  Any $\to$ third node & \$18--50K & More replicas or parallel training jobs. \\
  \bottomrule
\end{tabularx}}

\ReportSubsection[sec:dell]{Dell-configured pricing}
\TodoAnchor{aoa-dell}

Hardware would be ordered from Dell, configured to specification. Street prices
understate that cost because enterprise servers include vendor-qualified components,
warranty and support. A planning multiplier of 1.3 (range 1.15--1.4) is applied to
hardware until a quote is available; labor and operating costs are unchanged.

\TableTitle[tbl:dell]{Hardware and 3-year cost with Dell-configured pricing (USD)}
{\footnotesize
\begin{xltabular}{\textwidth}{@{}l Y{0.848} Y{1.091} Y{0.848} Y{1.213}@{}}
  \toprule
  Config. & Street-price hardware & Dell-configured hardware (range) & Dell hardware (expected) & 3-yr total, Dell-configured \\
  \midrule
  \endhead
  S1 & \$30--40K & \$35--56K & \$46K & \$81--121K \\
  S2 & \$39--50K & \$45--70K & \$57K & \$94--140K \\
  D1 & \$48--62K & \$55--87K & \$72K & \$117--174K \\
  \textbf{D3} & \$57--72K & \$66--101K & \textbf{\$83K} & \textbf{\$131--192K} \\
  S4 & \$58--73K & \$67--102K & \$85K & \$122--181K \\
  D4 & \$66--82K & \$76--115K & \$95K & \$144--211K \\
  \bottomrule
\end{xltabular}}
{\footnotesize 3-year totals add 3 years of operating cost, setup labor and upkeep (Table~\ref{tbl:tco}) to the Dell-configured hardware.}

\textbf{Notes for the Dell quote.}
\begin{itemize}
  \item Ask which current GPU-capable PowerEdge models (for example the 2U R-series and the
    GPU-dense XE-series) and Precision rack or tower workstations support the RTX PRO 6000
    Blackwell, how many GPUs each holds, and the power and cooling needed. The GPU count
    per chassis decides the growth path; confirm model numbers with the Dell account team.
  \item Dell servers typically use passive 600~W Server Edition GPUs instead of 300~W Max-Q
    cards; add about 0.3~kW per GPU to the power figures in the facility table.
  \item For Computer~1, request a storage-dense 2U PowerEdge with 12 or more 3.5-inch bays,
    redundant power, TPM~2.0 and iDRAC, with NVMe boot or VM drives.
  \item Dell support bundled in the quote (for example ProSupport) may replace part of the
    annual non-labor support cost assumed below.
  \item Request country-of-origin and TAA/Section~889 documentation with the quote.
\end{itemize}

\ReportSubsection[sec:sizing]{Sizing logic}
\TodoAnchor{aoa-bench}

\textbf{Weights memory} is parameters times bytes per parameter (FP16 = 2,
FP8 = 1, 4-bit $\approx$ 0.5), plus about 10\% overhead. The \textbf{KV cache}
per token is
\[
  \text{KV bytes/token} = 2 \times \text{layers} \times \text{KV heads} \times \text{head dim} \times \text{bytes per element}.
\]

\textbf{Example 1: gpt-oss-120b on one 96~GB GPU.} Approximate architecture: 36
layers, 8 KV heads, head dimension 64, with about half the layers using a short
sliding window, so only about 18 layers cache the full context. In BF16 the
cache is about $2 \times 18 \times 8 \times 64 \times 2~\text{B} \approx 37$~KB per
token, or about 4.7~GB for a 128K-token session. With
$96 \times 0.9 \approx 86$~GB and about 63~GB of weights, roughly 23~GB remains
for cache: about 5 simultaneous 128K sessions or 70 or more at 8K context.
Compute, not memory, usually limits short-prompt use, which is why one GPU is
rated for 3--8 simultaneous active users.

\textbf{Example 2: 70B dense model in FP8 on two GPUs (192~GB).} With 80 layers,
8 KV heads, head dimension 128 and an FP8 cache, the cost is
$2 \times 80 \times 8 \times 128 \times 1~\text{B} \approx 160$~KB per token (about
21~GB per 128K session). After about 70~GB of weights, about 103~GB remains:
roughly 5 full-128K sessions or 80 at 8K. A mixture-of-experts model is therefore
much more efficient than a dense 70B model for long-context, multi-user serving.

\[
  \text{Concurrent sessions} \approx \frac{\text{GPU memory} \times 0.9 - \text{weights}}{\text{average context tokens} \times \text{KV bytes per token}}
\]

\textbf{Throughput (order of magnitude).} Each GPU has about 1.8~TB/s memory
bandwidth. A small-active-parameter mixture-of-experts model can sustain roughly
1,000--3,000 aggregate output tokens/s per GPU under batching. Dense 70B models
are several times slower. \emph{Benchmark the actual model and real prompts
before purchase.}

\textbf{Fine-tuning memory (planning estimates).}

\TableTitle{Fine-tuning feasibility on 96 GB and 192 GB nodes}
{\small
\begin{xltabular}{\textwidth}{@{}Y{1.138}Y{0.828}Y{1.034}@{}}
  \toprule
  Method and model & Approx.\ GPU memory & Fits on \\
  \midrule
  \endhead
  Adapter (LoRA/QLoRA), 7--30B model & 20--70~GB & One 96~GB GPU \\
  QLoRA, 70B dense model (4-bit base) & 48--70~GB & One 96~GB GPU (tight at long sequence lengths) \\
  Adapter fine-tuning of gpt-oss-120b or similar & about 65--90~GB or more (verify with current tooling) & One 96~GB GPU is tight; two GPUs (192~GB) is comfortable \\
  Full fine-tuning (about 16 bytes per parameter with Adam in mixed precision) & 7B: about 112~GB; 30B: about 480~GB & 7B or smaller only on a 2-GPU node; larger is out of scope \\
  Pre-training a new LLM & Far beyond this budget & Not feasible \\
  \bottomrule
\end{xltabular}}

Adapter jobs on a few thousand to roughly one hundred thousand examples typically
run for hours to a few days on one GPU; confirm in the pilot. The GPUs connect
over PCIe (no NVLink), which is fine for serving and adapter training.

\ReportSubsection{Sizing check against the stated requirement}

Requirement: at least 20 named users, about 12 generating at the same time.
\begin{itemize}
  \item \textbf{Streaming load.} 12 users $\times$ 20 tokens/s is about 240 tokens/s,
    far below the roughly 1,000--3,000 tokens/s a mixture-of-experts model can sustain per
    GPU in aggregate.
  \item \textbf{Cache memory.} With gpt-oss-120b, 12 sessions of 16K tokens need about 7~GB;
    even 3 sessions at 128K (about 14~GB) plus 9 at 16K (about 5~GB) need about 19~GB, which
    fits in the roughly 23~GB left on one 96~GB GPU after the weights.
  \item \textbf{Limiting factor.} Processing many long prompts at once (prompt prefill) and
    keeping time to first token low. This is why one GPU is rated 3--8 simultaneous users;
    the rating is conservative and the pilot benchmark may show more.
  \item \textbf{Conclusion.} A 1-GPU serving node is marginal for 12 users. A 2-GPU node
    (rated 8--20) meets the need with margin, and Node~B can add capacity when it is not
    training. If named users grow well beyond 20, simultaneous use may exceed 20: add a GPU
    or use Node~B as a second serving node.
\end{itemize}

\ReportSubsection{Facility}
\TodoAnchor{aoa-facility}

\TableTitle{System power at load and circuit guidance}
{\small
\begin{tabularx}{\textwidth}{@{}l Y{0.6} Y{1.4}@{}}
  \toprule
  Config. & Power at load & Circuit and cooling \\
  \midrule
  S1 & 1.1--1.5~kW & One 208/240~V 20~A circuit. \\
  S2 & 1.6--2.1~kW & One 208/240~V 20~A circuit. \\
  D1 & 1.9--2.5~kW & One 208/240~V 20~A circuit (about 3.3~kW continuous). \\
  D3 & 2.4--3.1~kW & Two 20~A circuits, or one 30~A circuit, with separate PDUs per node. \\
  S4 & 2.3--3.1~kW & Two 20~A circuits or one 30~A circuit. \\
  D4 & 2.9--3.7~kW & Two 20~A circuits or one 30~A circuit; about 12,600~BTU/h at the top of the range. \\
  \bottomrule
\end{tabularx}}

Figures assume 300~W (Max-Q) GPUs: about 0.8--1.0~kW for N1, 1.3--1.6~kW for N2,
2.0--2.6~kW for N4 and 0.3--0.5~kW for Computer~1. If Dell supplies 600~W Server
Edition GPUs, add about 0.3~kW per GPU. The nodes are loud; place them
in a machine room or sound-attenuated rack.

\ReportSubsection[sec:procure]{Procurement considerations}
\begin{itemize}
  \item \textbf{Dell as vendor; TAA, Section 889 and NDAA compliance.} Dell is the stated
    vendor for configured-to-order servers. Request quotes for Computer~1 and for each
    candidate node configuration, with country-of-origin and TAA/Section~889
    documentation, and compare with the street-price planning figures in this report
    (Section~\ref{sec:dell}).
  \item \textbf{Memory and disk price volatility.} Server RAM and large HDD prices
    have been unstable. Quote early and consider buying less RAM now with expansion
    later.
  \item \textbf{Warranty and spares.} Three-year on-site coverage where available; hold
    a spare PSU and a spare HDD at minimum. Plan to retain failed drives (CUI)
    rather than returning them.
  \item \textbf{License servers.} Confirm with each software vendor that the license
    server may run in a virtual machine and what hardware binding applies, before
    finalizing the Computer~1 design.
  \item \textbf{Lead time and obsolescence.} Confirm GPU availability; plan for a
    3--5 year service life. The model layer can be replaced without hardware change.
\end{itemize}

% ===========================================================================
\ReportSection[sec:comparison]{Comparative Analysis}

\ReportSubsection[sec:tco]{Three-year cost}
\TodoAnchor{aoa-tco}

Hardware in this table is at street prices; Dell-configured totals are in
Table~\ref{tbl:dell}. Assumptions (also used in the Business Justification and accepted by
the sponsor for now): operator salary
\$100,000/year, or \$48.08/hour with no overhead multiplier. Upkeep: 2~hr/week
for the ML side (given) plus 1~hr/week for Computer~1 (assumed); a second node
adds 1~hr/week (assumed). Setup: 80--160~hr for Computer~1, 160--320~hr for the
first ML node, and 40--80~hr for a second node (all assumed).

\TableTitle[tbl:tco]{Estimated 3-year cost by configuration (USD)}
{\scriptsize
\begin{xltabular}{\textwidth}{@{}l Y{0.891} Y{1.366} Y{0.891} Y{0.95} Y{0.95} Y{0.952}@{}}
  \toprule
  Config. & Hardware & Setup labor & Annual non-labor opex & Annual upkeep & \textbf{3-yr total} & 3-yr ML only$^\dagger$ \\
  \midrule
  \endhead
  S1 & \$30--40K & 240--480~hr; \$11.5--23.1K & \$4--6.5K & 3~hr/wk; \$7.5K & \textbf{\$76--105K} & \$47--63K \\
  S2 & \$39--50K & 240--480~hr; \$11.5--23.1K & \$5--8K & 3~hr/wk; \$7.5K & \textbf{\$88--120K} & \$59--77K \\
  D1 & \$48--62K & 280--560~hr; \$13.5--26.9K & \$6--10K & 4~hr/wk; \$10K & \textbf{\$110--149K} & \$80--107K \\
  D3 & \$57--72K & 280--560~hr; \$13.5--26.9K & \$7--11.5K & 4~hr/wk; \$10K & \textbf{\$122--163K} & \$92--121K \\
  S4 & \$58--73K & 240--480~hr; \$11.5--23.1K & \$7--11K & 3~hr/wk; \$7.5K & \textbf{\$113--152K} & \$84--109K \\
  D4 & \$66--82K & 280--560~hr; \$13.5--26.9K & \$8--13K & 4~hr/wk; \$10K & \textbf{\$134--178K} & \$104--136K \\
  \bottomrule
\end{xltabular}}
{\footnotesize $^\dagger$ Excludes Computer~1 (about \$29--42K over 3 years, including its hardware, setup, operating cost and 1~hr/wk upkeep), which is needed for storage and licensing under any LLM alternative except the status quo. Non-labor opex covers support, software, power and cooling. 3-year total = hardware + setup labor + 3 $\times$ (non-labor opex + upkeep).}

\ReportSubsection{Scoring}

\TableTitle{Weighted scores (1 = poor, 5 = excellent; illustrative, revise with stakeholders)}
{\small
\begin{xltabular}{\textwidth}{@{}Y{1.0} r r r r r r@{}}
  \toprule
  Criterion (weight) & Alt 0 & Alt 2 & S2 & D1 & D3 & D4 \\
  \midrule
  \endhead
  CUI compliance and control (25\%) & 4 & 3 & 5 & 5 & 5 & 5 \\
  Capability and context (15\%) & 1 & 5 & 4 & 3 & 4 & 4 \\
  Concurrency and scalability (15\%) & 2 & 5 & 2 & 2 & 3 & 4 \\
  Resilience and continuity (15\%) & 3 & 5 & 2 & 4 & 4 & 4 \\
  3-year cost (15\%) & 5 & 3 & 4 & 3 & 3 & 2 \\
  Time to capability (5\%) & 5 & 3 & 3 & 3 & 3 & 3 \\
  Integration and customization (10\%) & 1 & 3 & 5 & 5 & 5 & 5 \\
  \midrule
  \textbf{Weighted score} & \textbf{3.00} & \textbf{3.90} & \textbf{3.70} & \textbf{3.70} & \textbf{4.00} & \textbf{4.00} \\
  \bottomrule
\end{xltabular}}

Interpretation: D3 scores as high as D4 at about \$9K less hardware at street prices, and it keeps a
failover path and a dedicated training node. The compliant cloud scores well on the
LLM alone but is sensitive to usage growth and to CUI contract and risk acceptance,
and it does not replace Computer~1. The scores are judgments, not measurements. D1 appears
in the table but does not meet the stated 12 simultaneous users.

\ReportSubsection{Break-even versus a compliant cloud LLM}
\[
  \text{Break-even (months)} = \frac{\text{On-premises ML upfront cost}}{\text{Monthly cloud cost} - \text{Monthly on-premises ML run cost}}
\]
Computer~1 is excluded because it is needed in either case. For D3 at Dell-configured
prices, the ML upfront cost is about \$64K of hardware (street \$49K $\times$ 1.3) plus
about \$14K of setup labor, or about \$78K, and the ML run cost is about \$1.25K/month
(non-labor opex plus upkeep). If cloud spend reaches \$4K/month, break-even is about
28 months. Averaged over 3 years the on-premises ML side costs about \$3.4K/month, so the
cloud is cheaper if spend stays below that. For 20--30 users a compliant cloud LLM is
therefore plausibly cheaper than D3 on usage cost alone unless usage is heavy. The case
for on-premises rests on keeping CUI inside the boundary, training on the NAS archive
without sending it off-site, no per-token cost growth, and the fact that Computer~1 is
needed anyway. Usage profiling in the pilot should drive this choice.

\ReportSubsection{Capability notes}
\begin{itemize}
  \item Open-weight models in 2026 handle most enterprise tasks (summarization,
    extraction, drafting, scripting) well but may trail the best proprietary models on
    the hardest reasoning and agentic tasks.
  \item At this budget the largest practical models are about 120B parameters
    (nodes with 96--192~GB) and a 4-bit 400B-class mixture of experts on a 4-GPU
    node. Quantization lowers quality slightly; validate against real workflows.
  \item Long-context quality degrades with length; use retrieval and chunking instead of
    relying on raw 128K-token prompts for everything.
\end{itemize}

% ===========================================================================
\ReportSection{Risk Assessment}

\TableTitle{Risks (L = likelihood, I = impact: H/M/L)}
{\footnotesize
\begin{xltabular}{\textwidth}{@{}Y{0.839} c c Y{1.161}@{}}
  \toprule
  Risk & L & I & Mitigation \\
  \midrule
  \endhead
  Computer 1 is a single point of failure for the file share, license servers and the LLM application tier & M & H & RAID-Z2 with a hot spare, redundant PSUs, UPS, tested VM restore; second gateway on Node~B. No second server (sponsor decision). \\
  RAID and snapshots share one machine with the data, so they do not protect against server loss, ransomware or site loss & M & H & Encrypted offline backup set rotated off-site (about \$2--4K), restore tested. \\
  A fine-tuned model states wrong clause numbers or requirements with confidence & M & H & Retrieval with citations as the primary mechanism; evaluate adapters against known answers; human review of any use in deliverables. \\
  Standards (for example NFPA~780) are copyrighted and the license may not allow AI indexing or training & M & M & Check license terms before ingesting; keep edition records. \\
  A model trained on restricted folders is used by people without access & L & H & Per-group adapters; access-aware retrieval; train only on material open to all users of a model. \\
  Autodesk or Bluebeam needs internet sign-in or cloud services outside the boundary & M & M & Verify licensing model and self-hosted options with the vendors before design freeze. \\
  License vendor restricts virtualization or binds licenses to hardware & M & M & Confirm terms before design freeze; use a dedicated host if required. \\
  Hardware must be TAA-compliant, raising cost 15--40\% & H & M & Confirm sourcing rules early; obtain OEM quotes alongside white-box. \\
  Authorization (ATO) delays, now covering file share and private cloud as well as the LLM & M & H & Engage the ISSO/ISSM early; one boundary and SSP for both computers. \\
  Fine-tuned models or embeddings inherit CUI sensitivity & M & M & Registry access control; include in backup and sanitization rules. \\
  Training jobs interfere with serving (single-node configs) & M & M & Run off-hours or dedicate GPUs; prefer dual-node for heavy training. \\
  Upkeep exceeds 3--4 hr/week during growth or after incidents & M & M & Named owner and backup person; configuration as code; scheduled maintenance windows. \\
  Under-sizing as adoption grows & M & M & Quotas, queueing, instrumented pilot; add a GPU at 70\% sustained utilization. \\
  GPU, RAM or disk supply and price volatility & M & M & Early quotes; accept equivalent parts. \\
  Model license or provenance concern & L--M & H & Legal review; approved-model list. \\
  Hallucination or misuse & H & M & Citations, human review, acceptable-use policy. \\
  Data leakage between users & L & H & Per-group access control, per-tenant indexes, red-team tests. \\
  Prompt injection through documents & M & M & Input sanitization; limited tool permissions; no autonomous external actions. \\
  Facility power, cooling and noise & M & M & Machine-room site check; 208/240~V circuits. \\
  \bottomrule
\end{xltabular}}

% ===========================================================================
\ReportSection{Implementation Roadmap}

\TableTitle{Phases}
{\small
\begin{xltabular}{\textwidth}{@{}Y{0.545}Y{1.636}Y{0.819}@{}}
  \toprule
  Phase & Activities & Output \\
  \midrule
  \endhead
  0. Discovery & Confirm user counts, workflows, storage and license inventory, sourcing rules (TAA), model policy, facility, recovery targets. & Validated requirements. \\
  1. Pilot and benchmark & Borrow or buy one ML node; benchmark 2--3 candidate models, an adapter fine-tune, and retrieval over a licensed standard such as NFPA~780 (non-sensitive) on representative prompts; test 12 simultaneous users; compare with the current tools. & Throughput, latency and quality data; configuration decision. \\
  2. Procurement and ATO preparation & Quotes, purchase order, SSP, STIG baselines, model-import procedure, license-vendor confirmations. & Approved design. \\
  3. Deploy Computer 1 and Node A & Install, harden, migrate file share and license servers, deploy gateway, retrieval and SSO; 20--40 pilot users. & Initial operating capability. \\
  4. Add Node B (and GPUs) & Dedicated training node, failover test, rolling-upgrade procedure; add GPUs as utilization triggers fire. & Full operating capability. \\
  5. Sustain & Quarterly model refresh, security review, restore test, usage reporting. & Continuous improvement. \\
  \bottomrule
\end{xltabular}}

\textbf{Expansion triggers:} sustained GPU utilization above 70\%, 95th-percentile time to first token above target,
or queue wait above 10~s. \textbf{Failover test:} run before production use and
after every major change.

% ===========================================================================
\ReportSection[sec:open]{Decision Requested and Open Items}
\TodoAnchor{aoa-open}

\textbf{Decision requested.} Approve configuration D3 (about \$64K of hardware at street
prices, about \$83K Dell-configured) plus integration effort to proceed with Phases 0--4.
Funding may be staged: Computer~1 and Node~A first, Node~B after the pilot confirms
demand and the training workload. If failover at the full 12 users is required, plan
Node~B with 2 GPUs (D4).

\textbf{Open items}
\begin{enumerate}
  \item Expected growth in named users beyond 20, and the workflow inventory by department.
  \item Autodesk and Bluebeam: which products, whether licenses are named-user or network,
    whether the enclave allows vendor sign-in, and which server components are needed
    (Section~\ref{sec:licenses}).
  \item Archive size, file formats (scanned PDFs, Office files, DWG/Revit, Bluebeam
    markups), folder and permission structure, and which documents to ingest first (for
    example NFPA~780, with license terms for AI use).
  \item Whether an offline backup copy is wanted in addition to disk redundancy inside
    Computer~1.
  \item Recovery targets (RTO/RPO) for the LLM service and for the license servers and file share.
  \item Type and size of training workloads (adapter fine-tuning, embeddings, small custom models).
  \item CUI categories in scope, including export-controlled data; contract restrictions.
  \item Dell quotes for Computer~1 and each node configuration; which models support the
    chosen GPUs; support level (Section~\ref{sec:dell}).
  \item Policy on open-weight models and model origin; whether cloud CUI processing is allowed.
  \item Facility readiness: power, cooling, noise, physical access.
  \item Who operates both computers and how much time they can dedicate.
  \item Approved marking text and handling rules for this document set.
  \item Funding profile and procurement vehicle; at least two vendor quotes per configuration.
  \item Server operating-system standard: Windows Server with Hyper-V and Active Directory,
    or Linux with KVM (affects Autodesk/Bluebeam hosting and file-share permissions).
  \item Failover target: is 3--8 users acceptable while Node~A is down, or must 12 be
    supported (Node~B with 2 GPUs)? Is Node~B expected to act as a second serving node?
  \item Which cloud fallback option(s) in Section~\ref{sec:fallback} to refine.
\end{enumerate}

% ===========================================================================
\ReportSection{Glossary and References}

\TableTitle{Terms used in this report}
{\footnotesize
\begin{xltabular}{\textwidth}{@{}Y{0.55} Y{1.45}@{}}
  \toprule
  Term & Meaning \\
  \midrule
  \endhead
  Adapter (LoRA, QLoRA) & A small set of extra weights trained on top of a fixed model to teach it new vocabulary, formats or tasks. QLoRA does this on a compressed (quantized) copy to save memory. \\
  Air-gapped & A network physically or logically isolated from the internet. \\
  ATO (Authority to Operate) & Formal approval to run a system, based on its security documentation and testing. \\
  CMMC & Cybersecurity Maturity Model Certification: a Department of Defense program requiring contractors to meet security practices such as NIST SP 800-171. \\
  Context window & The amount of text, counted in tokens, that the model can consider at once: the question, any supplied documents and the answer so far. \\
  CUI & Controlled Unclassified Information: unclassified government information that must be safeguarded. \\
  Dell-configured price & The price of a server built to order by Dell, including qualified parts, warranty and support. Here estimated as the street price $\times$ 1.3. \\
  DFARS 252.204-7012 & Defense acquisition clause requiring adequate security and incident reporting for covered defense information. \\
  Embedding & A list of numbers representing the meaning of a passage of text, used to search by similarity. \\
  Failover & Switching to a standby system, automatically or by hand, when the primary one fails. \\
  Fine-tuning & Further training of an existing model on selected examples to change how it behaves. \\
  FIPS 140-2/140-3 & Federal standards for validated cryptographic modules. \\
  GPU and VRAM & Graphics processing unit, and its on-board memory (VRAM), which holds the model and the conversation caches. \\
  HBA & Host bus adapter: the card that connects a set of drives to the server. \\
  Hypervisor and virtual machine (VM) & Software that runs several isolated virtual computers on one physical server. \\
  Inference & Running a trained model to produce an answer, as opposed to training it. \\
  KV cache & Memory the model keeps for each active conversation so it does not re-read the whole text for every new word; it grows with context length. \\
  LLM & Large language model. \\
  Mixture of experts (MoE) & A model design in which only a small part of the model is used for each word, making it faster than a dense model of the same size. \\
  Named and simultaneous users & People who have accounts, and people generating an answer at the same moment. \\
  NAS & Network-attached storage: shared file storage reached over the network. \\
  NIST SP 800-171 & NIST publication listing the security requirements for protecting CUI on non-federal systems. \\
  OCR & Optical character recognition: turning scanned pages into text. \\
  Open-weight model & A model whose trained weights are published for local use under a license. \\
  Prefill & The step where the model reads the prompt before it starts writing; very long prompts make it slower. \\
  Quantization & Storing model numbers with fewer bits (for example 4-bit) to save memory at a small cost in accuracy. \\
  RAID-Z2 & A disk arrangement that tolerates the failure of any two drives in the array. \\
  Reranker & A model that reorders retrieved passages by how well they answer the question. \\
  Retrieval (RAG) & Retrieval-augmented generation: finding relevant passages in an index of documents and giving them to the model so it answers from them, with citations. \\
  RTO and RPO & Recovery time objective (how long an outage may last) and recovery point objective (how much recent data may be lost). \\
  SIEM & Security information and event management: the central system that collects and analyzes security logs. \\
  SMB and NFS & Network file-sharing protocols used by Windows and by Linux/Unix systems. \\
  Snapshot & A read-only point-in-time copy of a file system that allows recovery from deletion or change. \\
  SSP & System Security Plan. \\
  STIG & Security Technical Implementation Guide: configuration hardening rules. \\
  TAA and Section 889 & Trade Agreements Act and NDAA Section~889: sourcing rules limiting country of origin and certain vendors. \\
  Tensor parallelism & Splitting one model across several GPUs so they work on each word together. \\
  Token & A piece of text, roughly three-quarters of a word, used as the unit for context size and speed. \\
  Vector database & A database that stores embeddings and finds the most similar ones quickly. \\
  \bottomrule
\end{xltabular}}

\textbf{Reference topics to cite in the final report (verify current versions):}
NIST SP 800-171 (Rev~2/3), 800-171A, 800-53; NIST AI RMF 1.0; 32~CFR Part~2002;
DFARS 252.204-7012; CMMC (32~CFR Part~170); NDAA Section~889; Trade Agreements Act;
vendor datasheets for the selected GPUs, servers and storage; model cards and licenses
for each candidate model; vLLM and SGLang documentation.


\end{document}
